% ---
% titre: Calcul flash: 9VP - théorie
% theme: NO
% chapitre: Nombres naturels et décimaux
% sous_chapitre: Calcul réfléchi et mental
% niveau: [9e]
% type: TH
% tags: [c-addition,c-criteres-divisibilite,c-division,c-decomposition-facteurs-premiers,c-droite-numerique,c-livrets, c-multiplication,c-nombres-decimaux,c-operations,c-priorite-operations,c-nombres-premiers,c-puissance,c-ppmc,c-soustraction, f-individuel,p-flash-calcul]

% difficulte: 1
% source: latex
% ---

\semaineinvisible
\header{Calcul mental: 9VP - Semaine 1 - Tables \& livrets} 

 \vfill
%-----------------------------------------------------------------------
\subsection{Les tables de multiplication}
 
\begin{regle}{Définition}
La table de $n$ liste tous les multiples de $n$ : $n \cdot 1, n \cdot 2, \ldots, n \cdot 12$.
\end{regle}
 
\vspace{6pt}
 \textbf{Tables à connaître par coeur}:
 
\vspace{4pt}
\def\arraystretch{1.04}
\begin{tabularx}{\linewidth}{|*{16}{>{\centering\arraybackslash}X|}}
\hline
\cellcolor{themeNO!50}$\cdot$ &\cellcolor{themeNO!30}1 & \cellcolor{themeNO!30}2 & \cellcolor{themeNO!30}3 & \cellcolor{themeNO!30}4 &\cellcolor{themeNO!30}5 & \cellcolor{themeNO!30}6 & \cellcolor{themeNO!30}7 & \cellcolor{themeNO!30}8 &\cellcolor{themeNO!30}9 & \cellcolor{themeNO!30}10 & \cellcolor{themeNO!30}11 & \cellcolor{themeNO!30}12 & \cellcolor{themeNO!30}13  & \cellcolor{themeNO!30}14  & \cellcolor{themeNO!30}15  \\
\hline
\cellcolor{themeNO!30}1  & 1 & 2 & 3 & 4 & 5 & 6 & 7 & 8 & 9 & 10 & 11 & 12 & 13 & 14 & 15 \\\hline
\cellcolor{themeNO!30}2  & 2 & 4 & 6 & 8 & 10 & 12 & 14 & 16 & 18 & 20 & 22 & 24 & 26 & 28 & 30 \\\hline
\cellcolor{themeNO!30}3  & 3 & 6 & 9 & 12 & 15 & 18 & 21 & 24 & 27 & 30 & 33 & 36 & 39 & 42 & 45 \\\hline
\cellcolor{themeNO!30}4  & 4 & 8 & 12 & 16 & 20 & 24 & 28 & 32 & 36 & 40 & 44 & 48 & 52 & 56 & 60 \\\hline
\cellcolor{themeNO!30}5  & 5 & 10 & 15 & 20 & 25 & 30 & 35 & 40 & 45 & 50 & 55 & 60 & 65 & 70 & 75 \\\hline
\cellcolor{themeNO!30}6  & 6 & 12 & 18 & 24 & 30 & 36 & 42 & 48 & 54 & 60 & 66 & 72 & 78 & 84 & 90 \\\hline
\cellcolor{themeNO!30}7  & 7 & 14 & 21 & 28 & 35 & 42 & 49 & 56 & 63 & 70 & 77 & 84 & 91 & 98 & 105 \\\hline
\cellcolor{themeNO!30}8  & 8 & 16 & 24 & 32 & 40 & 48 & 56 & 64 & 72 & 80 & 88 & 96 & 104 & 112 & 120 \\\hline
\cellcolor{themeNO!30}9  & 9 & 18 & 27 & 36 & 45 & 54 & 63 & 72 & 81 & 90 & 99 & 108 & 117 & 126 & 135 \\\hline
\cellcolor{themeNO!30}10  & 10 & 20 & 30 & 40 & 50 & 60 & 70 & 80 & 90 & 100 & 110 & 120 & 130 &140 & 150 \\\hline
\cellcolor{themeNO!30}11 & 11 & 22 & 33 & 44 & 55 & 66 & 77 & 88 & 99 & 110 & 121 & 132 & 143 & 154 &165 \\\hline
\cellcolor{themeNO!30}12 & 12 & 24 & 36 & 48 & 60 & 72 & 84 & 96 & 108 & 120 & 132 & 144 & 156 & 168 & 180 \\\hline
\cellcolor{themeNO!30}13 & 13 & 26 & 39 & 52 & 65 & 78 & 91 & 104 & 117 & 130 & 143 & 156 & 169 & 182 & 195 \\\hline
\cellcolor{themeNO!30}14 & 14 & 28 & 42 & 56 & 70 & 84 & 98 & 112 & 126 & 140 & 154 & 168 & 182 & 196 & 210 \\\hline
\cellcolor{themeNO!30}15 & 15 & 30 & 45 & 60 & 75 & 90 & 105 & 120 & 135 & 150 & 165 & 180 & 195 & 210 & 225 \\\hline
\end{tabularx}
\def\arraystretch{1}

\begin{astuce}{Stratégie}
La multiplication est commutative ($a \cdot b = b \cdot a$). On peut toujours retourner une multiplication.\\
\textit{Exemple :} Si on oublie $8 \cdot 7$, on calcule $7 \cdot 8 = 56$.
\end{astuce}

%-----------------------------------------------------------------------
\vspace{10pt}
\subsection{Les divisions}
 
\begin{regle}{Définition}
Si $a \cdot b = c$, alors $c : a = b$ et $c : b = a$.
\end{regle}
 
\vspace{6pt}
 
\begin{astuce}{Stratégie}
Pour calculer $63 : 9$, on se demande : \textit{\og $9$ fois combien font $63$ ?\fg{}} et on cherche dans la table de $9$ : $9 \cdot 7 = 63$, donc $63 : 9 = 7$.
\end{astuce}

\newpage
%---------------------------------------------------------------------------------------------------------------------------------------------------------------------------------------------------------------------
\semaineinvisible
\header{Calcul mental: 9VP - Semaine 2 - $\times$ et $\div$ par 10, 100, 1000} 

\vfill
\subsection{Les valeurs de positions}

{\footnotesize
\def\arraystretch{1.1}
\begin{tabularx}{\linewidth}{|*{14}{C|}}
\hline
\rotatebox{90}{Unité des millions}&\rotatebox{90}{Centaines des milliers\quad}&\rotatebox{90}{Dizaines des milliers}&\rotatebox{90}{Unité des milliers}&\rotatebox{90}{Centaine}&\rotatebox{90}{Dizaine}&\cellcolor{themeNO!30}\rotatebox{90}{Unité}&\rotatebox{90}{Dixième}&\rotatebox{90}{Centième}&\rotatebox{90}{Millième}\\\hline
$\cdot 1\,000\,000$&$\cdot 100\,000$&$\cdot 10\,000$&$\cdot 1\,000$&$\cdot 100$&$\cdot 10$&\cellcolor{themeNO!30}$\cdot 1$&$\cdot 0,1$&$\cdot 0,01$&$\cdot 0,001$\\\hline
\end{tabularx}
\def\arraystretch{1}}

\begin{regle}{Définitions}
\textcolor{themeColor}{Multiplier} par 10, c'est donner à chaque chiffre une valeur 10 fois plus grande. Le chiffre des unités devient donc le chiffre des dizaines, le chiffre des dixièmes devient celui des unités, etc.\\[7pt]
\textcolor{themeColor}{Diviser} par 10, c'est donner à chaque chiffre une valeur 10 fois plus petite. Le chiffre des unités devient donc le chiffre des dixièmes, le chiffre des dixièmes devient celui des centièmes, etc.
\end{regle}

\vspace{10pt}\begin{astuce}{Stratégie}

\def\arraystretch{1}
\vspace{-10pt}\begin{tabularx}{\linewidth}{rX|rX}
\multicolumn{2}{C|}{$\cdot$ (chaque chiffre avance)}& \multicolumn{2}{C}{$:$ (chaque chiffre recule) }\\\hline
$\cdot 10$ & 1 rang vers la gauche & $: 10$ &1 rang vers la droite \\
$\cdot 100$ & 2 rangs vers la gauche & $: 100$ & 2 rangs vers la droite \\
$\cdot 1000$ & 3 rangs vers la gauche & $: 1000$ & 3 rangs vers la droite \\
\end{tabularx}
\def\arraystretch{1}
\\[8pt]
Astuce: le nombre de zéros indique de combien de rangs chaque chiffre se déplace.
\end{astuce}


%-----------------------------------------------------------------------
\vspace{10pt}
\subsection{Exemples}
\def\arraystretch{1.05}
\begin{tabularx}{\linewidth}{|X|X|X|}
\hline
\cellcolor{themeNO!30}\textbf{Opération} & \cellcolor{themeNO!30}\textbf{Déplacement} & \cellcolor{themeNO!30}\textbf{Résultat} \\\hline
$3,7 \cdot 10$ & 1 rang vers la gauche & $37$ \\\hline
$3,7 \cdot 100$ & 2 rangs vers la gauche  & $370$ \\\hline
$3,7 \cdot 1000$ & 3 rangs vers la gauche  & $3700$ \\\hline
$3,7 : 10$ & 1 rang vers la droite & $0,37$ \\\hline
$3,7 : 100$ & 2 rangs vers la droite & $0,037$ \\\hline
$3,7 : 1000$ & 3 rangs vers la droite & $0,0037$ \\\hline
\end{tabularx}
\def\arraystretch{1}

\begin{astuce}{Rappel}
La virgule ne bouge jamais. Elle marque la frontière entre les nombres entiers et décimaux.
\end{astuce}

\newpage
%---------------------------------------------------------------------------------------------------------------------------------------------------------------------------------------------------------------------
\semaineinvisible
\header{Calcul mental: 9VP - Semaine 3 - Doubles \& moitiés}

\vfill
%-----------------------------------------------------------------------
\subsection{Doubles et moitiés}

\begin{regle}{Définitions}
Prendre \textcolor{themeColor}{le double} d'un nombre $n$, c'est le multiplier par 2. \hspace{10pt} $\Longrightarrow\quad n \cdot 2$\\[3pt]
Prendre \textcolor{themeColor}{la moitié} d'un nombre $n$, c'est le diviser par 2. \hspace{29pt} $\Longrightarrow\quad n : 2$
\end{regle}

\def\arraystretch{1.3}
\begin{tabularx}{\linewidth}{|X|X|}
\hline
\cellcolor{themeNO!30}\textbf{Expression} & \cellcolor{themeNO!30}\textbf{Calcul} \\\hline
Le double de 15,7 & $15,7 \cdot 2 = 31,4$ \\\hline
La moitié de 32,6& $32,6 : 2 = 16,3$ \\\hline
\end{tabularx}


\vspace{6pt}
\begin{astuce}{Stratégie}
Pour doubler un nombre décimal, on double chaque partie séparément.\\
\textit{Exemple :} $2 \times 3,4 = 2 \times 3 + 2 \times 0,4 = 6 + 0,8 = 6,8$
\end{astuce}

%-----------------------------------------------------------------------
\vspace{10pt}
\subsection{Multiplier et diviser par 0,5}

\begin{regle}{Définitions}
\textcolor{themeColor}{Multiplier} un nombre $n$ par $0,5$, c'est le diviser par 2. \hspace{11pt} $\Longrightarrow\quad n\cdot 0,5 = n : 2$\\[3pt]
\textcolor{themeColor}{Diviser} un nombre $n$ par $0,5$, c'est le multiplier par 2. \hspace{10pt} $\Longrightarrow\quad n : 0,5 = n \cdot 2$
\end{regle}

\begin{tabularx}{\linewidth}{|X|X|}
\hline
\cellcolor{themeNO!30}\textbf{Expression} & \cellcolor{themeNO!30}\textbf{Calcul} \\\hline
$48 \cdot 0,5$ & $48 : 2 = 24$ \\\hline
$124 : 0,5$ & $124 \times 2 = 248$ \\\hline
\end{tabularx}

%-----------------------------------------------------------------------
\vspace{10pt}
\subsection{Multiplier et diviser par 4}

\begin{regle}{Définitions}
\textcolor{themeColor}{Multiplier} un nombre $n$ par $4$, c'est le doubler deux fois de suite. \hspace{33pt} $\Longrightarrow\quad n\cdot 4 = n \cdot 2 \cdot 2$\\[3pt]
\textcolor{themeColor}{Diviser} un nombre $n$ par $4$, c'est prendre la moitié deux fois de suite. \hspace{10pt} $\Longrightarrow\quad n : 4 = n : 2 : 2$
\end{regle}

\begin{tabularx}{\linewidth}{|X|X|}
\hline
\cellcolor{themeNO!30}\textbf{Expression} & \cellcolor{themeNO!30}\textbf{Calcul} \\\hline
$13,1 \cdot 4$ & $13,1 \cdot 2 = 26,2$ \quad puis \quad $26,2 \cdot 2 = 52,4$ \\\hline
$46 : 4$ & $46 : 2 = 23$ \quad puis \quad $23:2=11,5$ \\\hline
\end{tabularx}
\def\arraystretch{1}
  
\newpage
%---------------------------------------------------------------------------------------------------------------------------------------------------------------------------------------------------------------------
\semaineinvisible
\header{Calcul mental: 9VP - Semaine 4 - Astuces $\times$5, $\times$9, $\times$11, $\times$25}

\vfill
%-----------------------------------------------------------------------
\subsection{Multiplier et diviser par 5}

\vspace{-4pt}\begin{regle}{Définitions}
Pour multiplier par 5, on multiplie par 10, puis on divise le résultat par 2.\\[2pt]
Pour diviser par 5,  on divise par 10, puis on multiplie le résultat par 2.
\end{regle}

\begin{tabularx}{\linewidth}{|p{3cm}|X|p{3cm}|}
\hline
\cellcolor{themeNO!30}\textbf{Calcul} & \cellcolor{themeNO!30}\textbf{Étapes} & \cellcolor{themeNO!30}\textbf{Résultat} \\\hline
$34 \cdot 5$ & $34 \cdot 10 = 340$ \quad puis \quad  $340 : 2=170$ & $170$ \\\hline
$70 : 5$ & $70 \cdot 2 = 140$ \quad puis \quad  $140 : 10=14$ & $14$ \\\hline
\end{tabularx}

%-----------------------------------------------------------------------
\subsection{Multiplier par 9 et par 11}

\vspace{-4pt}\begin{regle}{Définitions}
Pour multiplier par 9, on multiplie par 10, puis on soustrait le nombre de départ.\\[2pt]
Pour multiplier par 11, on multiplie par 10, puis on ajoute le nombre de départ.
\end{regle}

\begin{tabularx}{\linewidth}{|p{3cm}|X|p{3cm}|}
\hline
\cellcolor{themeNO!30}\textbf{Calcul} & \cellcolor{themeNO!30}\textbf{Étapes} & \cellcolor{themeNO!30}\textbf{Résultat} \\\hline
$25 \cdot 9$ & $25 \cdot 10 = 250$ \quad puis \quad  $250-25=225$ & $225$ \\\hline
$43\cdot 11$ & $43 \cdot 10 = 430$ \quad puis \quad  $430+43=473$ & $473$ \\\hline
\end{tabularx}

%-----------------------------------------------------------------------
\subsection{Multiplier et diviser par 25}

\vspace{-4pt}\begin{regle}{Définitions}
Pour multiplier par 25, on multiplie par 100, puis on divise le résultat par 4.\\[2pt]
Pour diviser par 25,  on divise par 100, puis on multiplie le résultat par 4.
\end{regle}

\begin{tabularx}{\linewidth}{|p{3cm}|X|p{3cm}|}
\hline
\cellcolor{themeNO!30}\textbf{Calcul} & \cellcolor{themeNO!30}\textbf{Étapes} & \cellcolor{themeNO!30}\textbf{Résultat} \\\hline
$48 \cdot 25$ & $48 \cdot 100 = 4800$ \quad puis \quad  $4800 : 4=1200$ & $1200$ \\\hline
$1600 : 25$ & $1600 \div 100 = 16$ \quad puis \quad  $16 \cdot 4=64$ & $64$ \\\hline
\end{tabularx}

%-----------------------------------------------------------------------
\subsection{Récapitulatif des astuces}

\vspace{-4pt}\begin{tabularx}{\linewidth}{|>{\centering\arraybackslash}p{2.5cm}|p{4cm}|X|}
\hline
\cellcolor{themeNO!30}\textbf{Opérations} & \cellcolor{themeNO!30}\textbf{Astuces} & \cellcolor{themeNO!30}\textbf{Exemples} \\\hline
$\cdot\,5$ & $\cdot\, 10$ puis $: 2$ & $67 \cdot 5 =  67 \cdot 10 : 2 = 670 : 2 = 335$ \\\hline
$:5$ & $: 10$ puis $\cdot\,2$ & $46 \div 5 = 46:10 \cdot 2 = 4,6 \cdot 2 = 9,2$ \\\hline
$\cdot\,9$ & $\cdot\, 10$ puis $- n$ & $88 \times 9 = 88 \cdot 10 - 88 = 880 - 88 = 792$ \\\hline
$\cdot\,11$ & $\cdot \,10$ puis $+ n$ & $3,7 \times 11 = 3,7 \cdot 10 + 3,7 = 37 + 3,7 = 40,7$ \\\hline
$\cdot\,25$ & $\cdot\, 100$ puis $: 4$ & $32 \times 25 = 32 \cdot 100 : 4 = 3200 :4 = 800$ \\\hline
$:25$ & $: 100$ puis $\cdot\, 4$ & $450 : 25 = 450 : 100 \cdot 4 = 4,5 \cdot 4 = 18$ \\\hline
\end{tabularx}

\newpage
%---------------------------------------------------------------------------------------------------------------------------------------------------------------------------------------------------------------------
\semaineinvisible
\header{Calcul mental: 9VP - Semaine 5 - Priorité des opération}

\vfill
%-----------------------------------------------------------------------
\subsection{La règle de priorité : PEMDAS}

\vspace{-4pt}\begin{regle}{Définition}
Pour calculer une expression, on suit toujours le même ordre, résumé par l'acronyme \textcolor{themeColor}{\textbf{PEMDAS}}.
\end{regle}

\def\arraystretch{1.2}
\begin{tabularx}{\linewidth}{|>{\centering\arraybackslash}p{3cm}|X|}
\hline
\cellcolor{themeNO!30}\textbf{Lettre} & \cellcolor{themeNO!30}\textbf{Signification} \\\hline
\textbf{P} & Parenthèses \\\hline
\textbf{E} & Exposants (puissances) \\\hline
\textbf{M} / \textbf{D} & Multiplications et divisions \\\hline
\textbf{A} / \textbf{S} & Additions et soustractions \\\hline
\end{tabularx}
\def\arraystretch{1}

\begin{astuce}{Rappel}
Si plusieurs opérations de même priorité se suivent, on calcule \textbf{de gauche à droite}.
\end{astuce}

%-----------------------------------------------------------------------
\subsection{Stratégie de calcul}

\vspace{-4pt}\begin{astuce}{Méthode PEMDAS}
\textbf{1.} Calculer d'abord ce qui est entre parenthèses.\\
\textbf{2.} Calculer les exposants.\\
\textbf{3.} Calculer les $\times$ et $\div$, de gauche à droite.\\
\textbf{4.} Calculer les $+$ et $-$, de gauche à droite.
\end{astuce}

%-----------------------------------------------------------------------
\subsection{Exemples}

\vspace{-4pt}\def\arraystretch{1.15}
\begin{tabularx}{\linewidth}{|>{\centering\arraybackslash}p{4cm}|X|}
\hline
\cellcolor{themeNO!30}\textbf{Expression} & \cellcolor{themeNO!30}\textbf{Calcul} \\\hline
$(3 + 4) \cdot 2$ & $(3 + 4) \cdot 2 = 7 \cdot 2 = 14$ \\\hline
$3 + 4 \cdot 2$ & $3 + 4 \cdot 2 = 3 + 8 = 11$ \\\hline
$2 + 3^2$ & $2 + 3^2 = 2 + 9 = 11$ \\\hline
$12 : 4 \cdot 3$ & $12 : 4 \cdot 3 = 3 \cdot 3 = 9$ \\\hline
$(35:7\cdot2)^2 -30 + 2$ & $(35:7\cdot2)^2 -30 + 2 = (5\cdot2)^2 -30 + 2= 10^2 -30 + 2= 100 -30 + 2=70+2 =72$ \\\hline
\end{tabularx}
\def\arraystretch{1}

\begin{astuce}{Erreur fréquente}
La division a la même priorité que la multiplication. On ne peut pas les regrouper entre eux.\\[7pt]
$12 : 4 \cdot 3 \ne 12 : 12 = 1$ (FAUX, de droite à gauche)\\[3pt]
$12 : 4 \cdot 3 = 3 \cdot 3 =9$ (CORRECT, de gauche à droite)
\end{astuce}

\newpage
\semaineinvisible
\header{Calcul mental: 9VP - Semaine 8 - Décimaux : droite graduée} 

\vfill
%-----------------------------------------------------------------------
\subsection{Placer un nombre décimal}

\begin{regle}{Définition}
Sur une droite graduée, chaque graduation représente un même écart. On repère d'abord entre quels nombres entiers se situe le décimal, puis on affine selon le pas de graduation.
\end{regle}

\vspace{6pt}
\begin{astuce}{Stratégie}
\textit{Exemple :} pour placer $2,4$ entre $2$ et $3$ (graduation par dixièmes), on compte $4$ petits traits après $2$.
\end{astuce}

%-----------------------------------------------------------------------
\vspace{10pt}
\subsection{Comparer deux décimaux}

\begin{regle}{Définition}
On compare d'abord la partie entière. Si elle est égale, on compare les dixièmes, puis les centièmes, etc.
\end{regle}

\vspace{6pt}
\begin{astuce}{Stratégie}
\textit{Exemple :} $3,45 < 3,5$ car $3,45 = 3,450$ et $450 < 500$ millièmes.
\end{astuce}

%-----------------------------------------------------------------------
\vspace{10pt}
\subsection{Encadrer un décimal}

\begin{regle}{Définition}
Encadrer un nombre, c'est trouver les deux multiples les plus proches (d'une unité, d'un dixième, etc.) entre lesquels il se situe.
\end{regle}

\vspace{6pt}
\def\arraystretch{1.2}
\begin{tabularx}{\linewidth}{|X|X|}
\hline
\cellcolor{themeNO!30}\textbf{Nombre} & \cellcolor{themeNO!30}\textbf{Encadrement} \\\hline
$7,3$ (à l'unité) & $7 < 7,3 < 8$ \\\hline
$7,36$ (au dixième) & $7,3 < 7,36 < 7,4$ \\\hline
\end{tabularx}
\def\arraystretch{1}

\newpage
\semaineinvisible
\header{Calcul mental: 9VP - Semaine 9 - Compléments à 10, 100, 1000} 

\vfill
%-----------------------------------------------------------------------
\subsection{Le complément}

\begin{regle}{Définition}
Le \textcolor{themeColor}{complément} d'un nombre à 10, 100 ou 1000, c'est le nombre qu'il faut lui ajouter pour atteindre ce total.
\end{regle}

\vspace{6pt}
\begin{astuce}{Exemple}
Le complément de $37$ à $100$ est $63$, car $37 + 63 = 100$.
\end{astuce}

%-----------------------------------------------------------------------
%\vspace{6pt}
\subsection{Méthode}

\begin{astuce}{Stratégie}
On avance par étapes : jusqu'à la dizaine suivante, puis jusqu'à la centaine suivante, puis jusqu'au total. On additionne ensuite ce qu'il a fallu ajouter à chaque étape.
\end{astuce}

\vspace{6pt}
\begin{astuce}{Exemple guidé : complément de $347$ à $1000$}
\textbf{1.} De $347$ à $350$ : il manque $3$.\\
\textbf{2.} De $350$ à $400$ : il manque $50$.\\
\textbf{3.} De $400$ à $1000$ : il manque $600$.\\[3pt]
On additionne : $3 + 50 + 600 = 653$. Le complément de $347$ à $1000$ est donc $653$.
\end{astuce}

%-----------------------------------------------------------------------
%\vspace{6pt}
\subsection{Compléments décimaux}

\begin{astuce}{Stratégie}
Même méthode avec les décimaux : on complète d'abord jusqu'à l'entier suivant, puis jusqu'au total.\\[3pt]
\textit{Complément de $3,4$ à $10$ :} $3,4 \to 4$ (il manque $0,6$), puis $4 \to 10$ (il manque $6$). Donc $0,6 + 6 = 6,6$.
\end{astuce}

%-----------------------------------------------------------------------
%\vspace{6pt}
\subsection{Exemples}

\def\arraystretch{1.1}
\begin{tabularx}{\linewidth}{|X|X|}
\hline
\cellcolor{themeNO!30}\textbf{Calcul} & \cellcolor{themeNO!30}\textbf{Complément} \\\hline
Complément de $7$ à $10$ & $3$ \\\hline
Complément de $45$ à $100$ & $55$ \\\hline
Complément de $7,5$ à $10$ & $2,5$ \\\hline
\end{tabularx}
\def\arraystretch{1}


\newpage
\semaineinvisible
\header{Calcul mental: 9VP - Semaine 10 - Critères de divisibilité} 

\vfill
%-----------------------------------------------------------------------
\subsection{Qu'est-ce qu'un critère de divisibilité ?}

\begin{regle}{Définition}
Un critère de divisibilité permet de savoir si un nombre est divisible par un autre \textbf{sans faire la division}, en observant simplement ses chiffres.
\end{regle}

%-----------------------------------------------------------------------
\vspace{10pt}
\subsection{Les critères à connaître}

\def\arraystretch{1.6}
\begin{tabularx}{\linewidth}{|c|X|}
\hline
Divisible par?&Critère\\\hline
2&Le nombre est pair, il finit donc par 0,2,4,6,8).\\\hline
3&La somme des chiffres formant le nombre est divisible par 3.\\\hline
4&Le nombre formé par les deux derniers chiffres du nombre est divisible par 4.\\\hline
5&Le nombre finit par 0 ou 5.\\\hline
6&Le nombre est simultanément divisible par 2 et 3.\\\hline
9&La somme des chiffres formant le nombre est divisible par 9.\\\hline
10&Le nombre finit par 0.\\\hline
25&Le nombre finit par 00, 25, 50 ou 75.\\\hline
50&Le nombre finit par 00 ou 50.\\\hline
100&Le nombre finit par 00.\\\hline
\end{tabularx}
\def\arraystretch{1}

%-----------------------------------------------------------------------
\vspace{10pt}
\subsection{Combiner les critères}

\begin{astuce}{Stratégie}
Pour un critère combiné (comme $6$), on vérifie les deux critères séparément : le nombre doit remplir les deux conditions à la fois.
\end{astuce}


\newpage
\semaineinvisible
\header{Calcul mental: 9VP - Semaine 11 - Facteurs premiers, PGDC \& PPMC} 

\vfill
%-----------------------------------------------------------------------
\subsection{Les nombres premiers}

\begin{regle}{Définition}
Un nombre est \textcolor{themeColor}{premier} s'il a exactement deux diviseurs : $1$ et lui-même.
\end{regle}

\vspace{6pt}
\begin{astuce}{À connaître}
Les nombres premiers jusqu'à $50$ : $2, 3, 5, 7, 11, 13, 17, 19, 23, 29, 31, 37, 41, 43, 47$.
\end{astuce}

%-----------------------------------------------------------------------
\vspace{10pt}
\subsection{Décomposition en facteurs premiers}

\begin{regle}{Définition}
Décomposer un nombre, c'est l'écrire comme un produit de nombres premiers.
\end{regle}

\vspace{6pt}
\begin{astuce}{Exemple}
$60 = 2 \times 2 \times 3 \times 5 = 2^2 \times 3 \times 5$
\end{astuce}

%-----------------------------------------------------------------------
\vspace{10pt}
\subsection{PGDC et PPMC}

\begin{regle}{Définition}
Le \textcolor{themeColor}{PGDC} de deux nombres est le plus grand nombre qui les divise tous les deux.\\[3pt]
Le \textcolor{themeColor}{PPMC} de deux nombres est le plus petit nombre qui est un multiple des deux.
\end{regle}

\vspace{6pt}
\begin{astuce}{Exemple}
$12 = 2^2 \times 3$ \quad et \quad $18 = 2 \times 3^2$\\[3pt]
PGDC : on garde les facteurs \textbf{communs}, avec le plus \textbf{petit} exposant : PGDC(12;18) $=2 \times 3 = 6$.\\
PPMC : on garde \textbf{tous} les facteurs, avec le plus \textbf{grand} exposant : PPMC(12;18) $=2^2 \times 3^2 = 36$.
\end{astuce}